\section{H1 and the convention problem}

The original intuitive statement was that the retained-state weight should
be a coarse vacuum amplitude.  That formulation merges three different
objects.

\begin{enumerate}[label=(\roman*)]
\item A positive transfer amplitude \(T\) has a principal right eigenfunction
\(\psi\), satisfying \(T\psi=\lambda_0\psi\).
\item The vacuum Doob transform uses the ratio \(\psi_j/\psi_i\) to convert
\(T\) into a normalized state-transition kernel.
\item A forward rare-state exploration uses a Lyapunov cost weight \(V\).  In
the vacuum-derived construction the natural orientation is
\(V_i\propto\psi_i^{-1}\), not \(V_i\propto\psi_i\).
\end{enumerate}

There is also an independent optimization: for a fixed bad subkernel \(Q\),
its positive Perron vector minimizes
\begin{equation}
 \sup_i\frac{(QV)_i}{V_i}.
\end{equation}
This optimal defect weight need not be the reciprocal of the vacuum
eigenfunction of a larger transfer operator.

The sharpened H1 is therefore:

\begin{quote}
When retained-state dynamics is obtained by a Doob transform of a positive
vacuum transfer amplitude, inverse vacuum amplitude supplies a canonical
forward defect-cost weight.  For the resulting bad subkernel, Perron
optimization gives the best positive weighted ceiling.  Identifying either
vector with a physical Yang--Mills vacuum requires a reflection-positive
temporal transfer interpretation.
\end{quote}

This formulation separates physical time evolution, RG blocking, and an
exploration kernel.  A positive eigenvector of the latter two is not a
physical vacuum merely because the same linear algebra appears.
